% Introduction to the EU - Week 14 Cornell Final Notes (Spring 2026, Mon/Tue)
% Compile: pdflatex introduction-to-eu-final-cornell-master-notes-mon-tue-spring2026.tex
\documentclass[10pt,a4paper]{article}
\usepackage[utf8]{inputenc}
\usepackage[T1]{fontenc}
\usepackage[margin=1.2cm]{geometry}
\usepackage{xcolor}
\usepackage{array}
\usepackage{longtable}
\usepackage{enumitem}
\setlist[itemize]{leftmargin=1.2em,itemsep=1pt,topsep=2pt}
\definecolor{eublue}{RGB}{18,48,100}
\definecolor{soft}{RGB}{245,248,252}
\setlength{\parindent}{0pt}
\setlength{\parskip}{0.35em}

\newcommand{\cornell}[3]{
\subsection*{#1}
\begin{longtable}{|p{0.22\linewidth}|p{0.74\linewidth}|}
\hline
\rowcolor{soft}\textbf{Cue / Questions} & \textbf{Notes} \\
\hline
#2 & #3 \\
\hline
\end{longtable}
}

\begin{document}
\begin{center}
{\Large\bfseries\color{eublue} Introduction to the EU - Week 14 Final Cornell Master Notes}\\
{\small Spring 2026 (typical class days: Mon/Tue) | Consolidated from weekly class notes, reading summaries, and exam prep blocks}
\end{center}

\section*{Performance and Focus Rules}
\begin{itemize}
  \item \textbf{History/EU risk marker from your note:} 5.5 (C+) or higher = aprobado bajo (AP) baseline.
  \item \textbf{Global fail thresholds:} WIP 3.49--0; UCM 4.99--0; suspenso = 4.
  \item \textbf{Priority rule:} this course should receive extra finals focus because it is in your ``focus most'' group.
\end{itemize}

\section*{Cornell Notes by Exam Theme}

\cornell{1) Integration Logic and Institutional Architecture}
{How does integration deepen?\\
Which institutions do what?\\
Where is legitimacy located?}
{
\textbf{Core logic:} EU integration evolves through institutional bargaining, legal development, and policy spillovers.\\
\textbf{Parliament/Council/Commission triangle:} legislation emerges from inter-institution negotiation, not one-center command.\\
\textbf{European Council role:} strategic direction under high political salience.\\
\textbf{Exam strategy:} define institution -> mechanism -> consequence for policy output.\\
\textbf{Common mistake:} listing institutions without explaining interdependence.
}

\cornell{2) CJEU, Primacy, and Integration Through Law}
{Why does legal primacy matter?\\
How does judicial interpretation shape integration?\\
Where do tensions appear?}
{
\textbf{Primacy doctrine:} EU law precedence secures uniform application across member states.\\
\textbf{CJEU function:} legal coherence, treaty interpretation, and enforcement support.\\
\textbf{Political-legal tension:} national constitutional claims can challenge practical compliance.\\
\textbf{Case-answer structure:} principle -> conflict -> institutional response -> integration implication.\\
\textbf{High-value point:} law is not neutral background; it is a core integration engine.
}

\cornell{3) EMU, ECB, and Economic Governance}
{What is EMU solving?\\
What is ECB mandate?\\
How do monetary/fiscal tensions emerge?}
{
\textbf{EMU purpose:} monetary coordination and currency stability under shared institutional rules.\\
\textbf{ECB role:} monetary policy with price-stability orientation and systemic confidence management.\\
\textbf{Governance challenge:} asymmetric shocks test solidarity and adjustment mechanisms.\\
\textbf{Exam framing:} explain institutional competence boundaries before normative judgment.\\
\textbf{Frequent error:} mixing monetary policy tools with fiscal competence assumptions.
}

\cornell{4) Citizenship, Identity, and Enlargement}
{What rights define EU citizenship?\\
How does enlargement reshape identity?\\
Why is migration politically central?}
{
\textbf{Citizenship layer:} adds cross-border rights beyond national frameworks.\\
\textbf{Identity debate:} institutional belonging, political narratives, and solidarity are uneven across states.\\
\textbf{Enlargement effect:} governance complexity increases while integration scope expands.\\
\textbf{Migration relevance:} intersection of rights, security, and political legitimacy.\\
\textbf{Exam tip:} avoid binary framing; discuss both integration gains and governance stress.
}

\cornell{5) Rule of Law and Policy Futures (Green Deal, Governance Resilience)}
{What is rule-of-law backsliding?\\
How does EU respond?\\
What do future policy agendas require?}
{
\textbf{Rule-of-law baseline:} legality, judicial independence, and rights protection as system conditions.\\
\textbf{Backsliding risk:} internal divergence can weaken legal and political cohesion.\\
\textbf{Policy future:} climate and industrial transitions require coordinated governance capacity.\\
\textbf{Analytical sentence:} ``EU policy ambition depends on institutional legitimacy and enforceable compliance mechanisms.''\\
\textbf{Finals focus:} write balanced answers with mechanism-based explanations, not only normative claims.
}

\cornell{6) Chronological Roots (1900--1950) and Postwar ``Atmosphere''}
{Why does the professor stress early 20th-century sequence?\\
How did WWI and Versailles reshape Europe?\\
What linked economic crisis to political extremism?\\
What was unique about 1945--1950 Western choices?}
{
\textbf{1914--1918 (WWI):} massive destruction; collapse or weakening of empires (e.g., Austro-Hungarian, Ottoman, Russian).\\
\textbf{1919 (Versailles):} territorial and economic settlement; strains and resentment contributed to interwar instability (especially German context in standard narratives).\\
\textbf{1920s--1930s:} depression and polarization; rise of authoritarian regimes (e.g., Fascism in Italy; Nazism in Germany).\\
\textbf{1939--1945 (WWII):} peak destruction; foundational shock for later cooperation projects.\\
\textbf{1945--1950 ``atmosphere'':} Europe divided (Iron Curtain framing); Western leadership sought to embed peace partly by binding key economies---especially \textbf{France and Germany}---so war became materially harder.\\
\textbf{Exam link:} integration is historically situated; institutions did not appear from abstract theory alone.
}

\cornell{7) EU as ``Union of Unions'': Layers, Schengen, Enlargement, and Rules}
{What does ``union of previous unions'' mean in practice?\\
What is Schengen and who is in/out?\\
Who are the founding fathers and the first core?\\
What enlargement waves matter for exams?\\
What are the Copenhagen Criteria and Lisbon (2007)?\\
What are the ``two faces'' of Europe?}
{
\textbf{Complexity thesis:} the EU bundles overlapping projects---not one single ``club''---e.g., \textbf{Schengen} (borderless internal movement for participants), \textbf{Eurozone} (single currency subset), and \textbf{Customs Union} (common external tariff, free trade inside).\\
\textbf{Schengen Agreement:} landmark 1985 / effective border-removal logic associated with 1995 wave in standard teaching; today described as \textbf{29} participating states with passport-free internal crossings for routine travel.\\
\textbf{Schengen breadth:} includes some \textbf{non-EU} participants (e.g., Norway, Switzerland in common lecture examples).\\
\textbf{Exceptions (as taught):} Ireland and Cyprus highlighted as Schengen exceptions in lecture.\\
\textbf{Founding fathers / Schuman line:} figures such as \textbf{Robert Schuman} (France) framed coal and steel pooling so war among industrial cores was ``not only unthinkable but materially impossible.''\\
\textbf{1951 / Core Six:} France, West Germany, Italy, Benelux as launching nucleus of the community project.\\
\textbf{1973 first enlargement:} UK, Ireland, Denmark (classic exam anchor).\\
\textbf{2004 big bang:} post--Cold War eastern enlargement (e.g., Poland, Estonia, and other accession states in that wave).\\
\textbf{Copenhagen Criteria (membership eligibility):} (1) stable democracy and institutions; (2) rule of law, rights, minority protection; (3) functioning market economy able to cope with competitive pressure (EU terminology paraphrased for notes).\\
\textbf{Acquis communautaire (tie-in):} candidate states must adopt accumulated EU law---membership is legal incorporation, not only diplomacy.\\
\textbf{Lisbon Treaty (2007):} institutional modernization; stronger \textbf{European Parliament} role and claims to greater transparency/legibility for citizens.\\
\textbf{Two faces:} \textit{Internal Schengen face}---freedom of movement inside the area for those covered; \textit{External face}---harder border and admission logic toward non-members or perceived non-democratic neighbors (e.g., Mediterranean and eastern neighborhood in lecture examples).\\
\textbf{Democracy requirement:} membership presupposes democratic systems; helps explain why some rich non-EU states participate in parts of the acquis/Schengen space while authoritarian neighbors do not.\\
\textbf{Motto:} \textit{United in Diversity} (official EU motto, 2000).\\
\textbf{Supranationalism cue:} pooled sovereignty through rules enforced beyond pure intergovernmental veto---link to customs union and court-driven integration elsewhere in this file.\\
\textbf{Exam phrasing:} ``The EU is historically rooted in war prevention, legally layered through overlapping treaties, and politically contested at its external boundary.''
}

\cornell{8) Course Introduction, Assessment, Institutions, and Critical History}
{What does the syllabus promise you will learn?\\
How is power organized at EU level (three branches framing)?\\
How many member states and what identity tension?\\
What are the graded components and attendance rules?\\
How does the syllabus frame Copenhagen in three pillars?\\
What is the ``official'' vs ``realist'' story?}
{
\textbf{Course objective:} European history, contemporary crises, and institutional \textbf{complexity} of the Union (not a simple ``international organization'' shortcut).\\
\textbf{Historical frame:} post-WWII project to reduce war risk through economic interdependence; early core around coal and steel cooperation to limit mutual armament/tariff rivalry (Schuman/ECSC lineage in narrative).\\
\textbf{Founding treaty cue:} \textbf{Treaty of Rome (1957)} as foundational reference alongside Paris/ECSC threads in class.\\
\textbf{Three powers (intro framing):} \textbf{Executive}---European Commission; \textbf{Legislative}---European Parliament (direct elections every five years) with Council role developed elsewhere in course; \textbf{Judicial}---Court of Justice of the EU (CJEU/ECJ).\\
\textbf{Membership:} \textbf{27} member states after UK withdrawal framing; persistent tension between \textbf{national} and \textbf{European} identity claims.\\
\textbf{Assessment (verify on official syllabus):} class \textbf{presentation} (e.g., 25\% weighting in transcript---two texts, ~20--25 min analysis + ~15 min debate); \textbf{midterm} with short exercises on texts/images (e.g., Treaty of Paris, coal/steel); \textbf{final} with structured questions (e.g., four per unit) requiring preface + synthesis; \textbf{research essay} (~3,500 words) on approved topic (examples: sport, religion, internet policy in the EU).\\
\textbf{Attendance / punctuality (as stated in class):} two absences without penalty; third absence may cost points; lateness after ~10 minutes tracked.\\
\textbf{Copenhagen Criteria (syllabus grouping):} (1) \textbf{Political}---stable democratic institutions, rule of law, human rights; (2) \textbf{Economic}---functioning market economy; (3) \textbf{Legal}---acceptance of the \textbf{acquis communautaire} (full body of rights and obligations).\\
\textbf{Vocabulary / lenses:} \textbf{Supranationalism}---pooled competence beyond pure national veto; \textbf{customs union}---common external tariff, free circulation inside; \textbf{butterfly effect}---crisis in one member (e.g., Greece, Spain) spills across the system; \textbf{colonial/decolonial} lens---how empire and its aftermath shape EU external relationships; \textbf{blue-collar/white-collar shift}---linked to populist and EU-skeptic voting patterns in lecture framing.\\
\textbf{Official vs realist history:} ``founding fathers'' peace narrative is necessary pedagogically, but \textbf{realist} reading adds economic interests, power rivalries, and colonial legacies that explain ongoing ``tense moments.''\\
\textbf{Exam phrasing:} ``EU answers should pair institutional accuracy with reflexivity---mechanism plus critical context, not only heroic integration folklore.''
}

\cornell{9) Post-WWII Origins (1945--1955): Marshall, ECSC, NATO, and Failed EDC}
{What was the post-1945 problem atmosphere?\\
How did US aid condition European cooperation?\\
What did the Council of Europe add beyond economics?\\
What did Schuman propose and what was ECSC?\\
How did NATO relate to the German question?\\
Why did the European Defence Community fail?\\
What is the ``midwife'' metaphor?}
{
\textbf{Post-war setting:} physical and economic ruin; elite goal to make renewed total war \textbf{not only unthinkable but materially impossible} via interlocked industry and oversight.\\
\textbf{Marshall Plan (1947):} US reconstruction aid (European Recovery Program; lecture cites order of magnitude \textbf{\$13 billion}); required coordinated European input, spurring institutions such as \textbf{OEEC (1948)} in Paris.\\
\textbf{Council of Europe (1949):} Strasbourg---symbolic Franco-German midpoint; emphasized \textbf{political union} and \textbf{human rights} discourse, not only market integration.\\
\textbf{Schuman Declaration (1950):} French Foreign Minister \textbf{Robert Schuman} proposed pooling French and West German \textbf{coal and steel} under a joint \textbf{High Authority}---birth narrative of supranational economic control of war-making inputs.\\
\textbf{ECSC (1951):} European Coal and Steel Community; \textbf{Core Six}---France, West Germany, Italy, \textbf{Benelux} (Belgium, Netherlands, Luxembourg). First major transfer of national \textbf{sovereignty} to a \textbf{supranational} executive.\\
\textbf{NATO (1949):} Atlantic military alliance (US, Canada, Western Europe) against Soviet pressure; classic capsule: ``keep the Russians out, the Americans in, and the Germans down.''\\
\textbf{Cold War split:} \textbf{Iron Curtain} metaphor for Soviet bloc isolation versus Western reconstruction/integration lane.\\
\textbf{EDC (1952--54):} European Defence Community---plan for a merged European army; \textbf{failed} when the French Assembly withheld ratification, fearing national sovereignty loss and uncontrollable German rearmament.\\
\textbf{Contrast for exams:} economic pooling (ECSC) succeeded where military integration (EDC) triggered sharper sovereignty anxieties.\\
\textbf{``Midwife'' metaphor:} three channels delivering ``new Europe''---(1) \textbf{economic} resource management (coal, steel, Marshall money); (2) \textbf{political} rights and assembly (Council of Europe); (3) \textbf{military} collective defence (NATO) plus the \textbf{abortive} European army project.\\
\textbf{Political cartoon cue:} \textbf{Tower of Babel} trope mocking early unity (languages/cultures) as obstacle to single army or government.\\
\textbf{German question (realist core):} French distrust of German industrial revival; shared coal/steel supervision aimed to steer capacity toward peaceful goods while providing transparency---steel for cars, not tanks, under European board logic.\\
\textbf{Acquis reminder:} later membership requires accepting full \textbf{acquis communautaire}; here the ECSC inaugurates the supranational precedent.\\
\textbf{Exam phrasing:} ``1945--55 grafts reconstruction economics onto security dilemmas: integration starts where war industries become jointly governed because diplomacy alone could not settle the German power problem.''
}

\cornell{10) From EEC Rome (1957) to Empty Chair and CAP (1950s--1960s)}
{Why did military union fail while economic union advanced?\\
What did the Treaty of Rome create and promise?\\
How was the early EEC governed (four bodies)?\\
What was the Empty Chair crisis and Luxembourg Compromise?\\
How did CAP fit French and consumer politics?\\
How did EEC differ from EFTA?}
{
\textbf{EDC failure (deeper lecture emphasis):} European Defence Community collapsed when French politics refused a merged army---\textbf{sovereignty} fears linked to \textbf{de Gaulle} line and \textbf{French Communist} opposition in standard classroom narrative, alongside wider Assembly reluctance on German rearmament.\\
\textbf{``Founding Fathers'' pivot:} \textbf{Jean Monnet} and \textbf{Robert Schuman} route---when defence pooling proved toxic, shift to \textbf{economic integration} (coal/steel then broader market) as slower but more durable ``federalism by sector.''\\
\textbf{Treaty of Rome (1957):} established \textbf{EEC}---aim of a \textbf{common market} by dismantling internal barriers (lecture metaphor: thousands of kilometres of border frictions) and lifting post-war \textbf{living standards}.\\
\textbf{Early EEC institutions:} \textbf{Commission}---independent treaty ``supervisor'' / initiator (often tied to \textbf{technocracy} critique); \textbf{Council of Ministers}---national executives where much real power sat; \textbf{European Parliament / Assembly}---initially \textbf{consultative} only (weak co-legislator era); \textbf{Court of Justice}---uniform application of Community law across states.\\
\textbf{NATO recap:} separate security pillar---``keep the Americans in, the Russians out, and the Germans down''---parallel to but not identical with EEC legal integration.\\
\textbf{``Empty Chair'' crisis (1965):} de Gaulle withdrew French officials from the Council about six months to block \textbf{Qualified Majority Voting (QMV)} perceived to dilute French control.\\
\textbf{Luxembourg Compromise (1966):} political settlement allowing a member to insist on consensus when it claims a \textbf{vital national interest}---sometimes glossed as a \textbf{qualified veto} shadow over QMV in practice.\\
\textbf{Dual reality (exam thesis):} treaty rhetoric and court doctrine push ``ever closer union,'' while high-politics vetoes show large states still steering when interests collide.\\
\textbf{CAP (1962):} \textbf{Common Agricultural Policy}---price/income support framework associated with French farmer politics; raised food costs and budget share debates but embedded agricultural interests inside the Community.\\
\textbf{EEC (Six) vs EFTA (Seven):} \textbf{EEC}---Franco-German core, \textbf{supranational} law and institutional ambition toward broad economic (and latent political) union; \textbf{EFTA}---UK-led club emphasizing \textbf{intergovernmental} \textbf{free trade} while preserving national sovereignty (``two Europes'' competitive story in the 1960s).\\
\textbf{Vocabulary:} \textbf{QMV}---weighted majority decisions; \textbf{intergovernmentalism}---state-controlled bargaining (EFTA/NATO diplomatic logic) vs \textbf{supranationalism}---pooling sovereignty in bodies like Commission/Court.\\
\textbf{Exam phrasing:} ``1960s integration alternates Monnet-style supranational market-building with Gaullist sovereignty re-assertion; CAP and the Luxembourg Compromise show economics and national vetoes co-constituting today's EU.''
}

\cornell{11) Gaullist Vetoes, CAP, and the 1973 Oil Shock (1950s--1970s)}
{How did de Gaulle frame British membership and use vetoes?\\
What CAP side-effects appear in lecture beyond ``French farmers''?\\
What triggered the 1973 oil crisis and what was stagflation?\\
Why did the crisis expose EEC weakness?\\
How did states respond with energy policy?\\
What are the ``three realities'' of European institutions?\\
How does Luxembourg cast a long shadow?}
{
\textbf{Gaullist line:} President \textbf{Charles de Gaulle} treated \textbf{supranational} federal hints with suspicion and saw UK entry as possible \textbf{US ``Trojan horse''}; used veto twice to block British accession to the \textbf{EEC} before eventual UK membership later (timeline beyond this block).\\
\textbf{Rome/EEC anchor:} \textbf{Treaty of Rome (1957)} builds the \textbf{common market}---internal barrier removal (lecture: ~4{,}000 km metaphor), Commission supervision, Council of national ministers, Court---ties to \S 10.\\
\textbf{CAP costs:} beyond price support for French agriculture, lecture cites consumer burden and symbolic \textbf{``butter mountains''} / surplus stockpiles from overproduction incentives.\\
\textbf{1973 oil shock:} \textbf{OPEC} leverage after Middle East conflict (\textbf{Yom Kippur War} context); dramatic price spike (lecture: order of \textbf{400\%} increase framing) ends post-war ``golden age'' growth narrative.\\
\textbf{Stagflation:} simultaneous \textbf{stagnation} and \textbf{inflation} (with unemployment stress)---macro breakdown unfamiliar to 1950s--60s boom mindset.\\
\textbf{Crisis governance lesson:} member states often reacted \textbf{nationally} (energy rationing, industrial bailouts) rather than through a unified Community playbook---exposing limited \textbf{political union} capacity while economics were integrated on paper.\\
\textbf{Energy pivot examples:} France accelerated \textbf{nuclear} investment; broad \textbf{conservation} campaigns (lower thermostats, \textbf{55 mph} caps where applied, \textbf{daylight saving} tweaks) show crisis managerialism outside Brussels coordination.\\
\textbf{``Three realities'' (parallel pillars):} (1) \textbf{Council of Europe}---Strasbourg human-rights and plural democracy talk; (2) \textbf{Economic Communities (ECSC/EEC)}---market rules and \textbf{supranational} law; (3) \textbf{NATO}---hard security still heavily \textbf{US}-anchored, distinct from EEC legal order.\\
\textbf{Luxembourg Compromise legacy:} post--Empty Chair practice lets states claim \textbf{vital interest} blockages---lecture argues this \textbf{slowed} deeper integration for decades by privileging nationalist brake pedals over automatic majoritarian community interest.\\
\textbf{Vocabulary recap:} \textbf{Common market}---free movement of goods (and, in fuller Single Market telling, services/people/capital); \textbf{intergovernmentalism} vs \textbf{supranationalism} helps decode who decides under shock conditions.\\
\textbf{Exam phrasing:} ``The 1973 shock reveals a decoupling: integrated market expectations met intergovernmental crisis management, prefiguring later EMU and energy debates.''
}

\cornell{12) 1973 Oil Crisis Deep Dive: Dependence, Butterfly Effect, and Iberian Friction}
{How does the lecture date the ``Golden Age'' and its end?\\
How did OPEC ``weaponize'' oil politically?\\
What dependency statistic illustrates European vulnerability?\\
What broke beyond prices (social contract, employment)?\\
How did national energy strategies diverge (France vs UK)?\\
Why does Gibraltar still matter in EU--UK--Spain politics?\\
How does ``butterfly effect'' connect Middle East war to daily life?}
{
\textbf{Golden Age framing:} roughly \textbf{1950--1973} rapid growth and reconstruction confidence in Western Europe; trajectory breaks with the \textbf{1973 oil shock} tied to the \textbf{Yom Kippur War}---extends \S 11 chronology with clearer periodization.\\
\textbf{OPEC logic:} exporters coordinated supply/pricing power, using oil as leverage against states perceived as backing \textbf{Israel}---``oil weapon'' beyond pure market clearing.\\
\textbf{Structural dependence:} lecture highlights heavy European reliance on petroleum---e.g., ca.\ \textbf{65\%} of electricity linked to oil in the period framing cited in class---making the cartel shock instantly macro-critical.\\
\textbf{Stagflation \& sociology of crisis:} simultaneous inflation and slowdown unraveled post-war \textbf{social contract} expectations; energy input costs propagated into \textbf{food} and \textbf{manufacturing}, feeding a modern \textbf{unemployment} wave unfamiliar after boom years.\\
\textbf{Failed common EEC energy statecraft:} Community institutions did not produce a single crisis \textbf{energy union} response; capitals improvised---\textbf{France} doubled down on \textbf{nuclear}; \textbf{UK} pursued \textbf{North Sea} hydrocarbons alongside rationing tropes---parallel national industrial strategies.\\
\textbf{Conservation toolkit (recap):} public-building thermostat cuts, \textbf{55 mph} motor limits where adopted, \textbf{daylight saving} adjustments---micro-austerity as substitute for pooled foreign policy.\\
\textbf{Institutional drag:} \textbf{Luxembourg Compromise} veto culture linked to Gaullist \textbf{intergovernmental} reflex meant even petroleum diplomacy struggled to centralize---integration \textbf{momentum} cooled just when supranational coordination might have helped.\\
\textbf{Spain, Gibraltar, Ceuta/Melilla:} \textbf{Gibraltar} remains UK territory contested by \textbf{Spain}; \textbf{Ceuta} and \textbf{Melilla} are Spanish exclaves in North Africa---recurring sovereignty irritants shaping bilateral and EU--UK border politics in lecture examples.\\
\textbf{Butterfly effect:} distant Middle Eastern conflict propagates to \textbf{cold classrooms in Spain}, slower French traffic, and broader austerity gestures---illustrates \textbf{global resource dependency} and why \textbf{energy transition} stayed a headline issue into contemporary EU agendas (incl.\ 2026 policy talk in transcript).\\
\textbf{Exam phrasing:} ``1973 proves external shocks exploit internal integration gaps: without shared energy governance or foreign-policy fusion, market Europe still dances to geopolitical commodity politics.''
}

\cornell{13) Delors Era, Single European Act, Schengen, and Maastricht (1985--1993)}
{Who led the Commission and what enlargements fit his tenure?\\
What was ``southern enlargement'' and why structural/cohesion money?\\
What did the Single European Act achieve?\\
How does Schengen fit the 1990s free-movement story?\\
How did Thatcher represent UK pushback?\\
What did Maastricht (1992) create and what were the three pillars?\\
What are opt-outs and variable geometry?\\
What is ``backdoor federalism'' in lecture terms?}
{
\textbf{Jacques Delors (Commission President ca.\ \textbf{1985--1995}):} presidency associated with relaunching ``1992'' internal-market momentum, managing \textbf{1986 Iberian accession}, and navigating \textbf{German reunification} integration within EC/EU structures.\\
\textbf{Southern enlargement (1986):} \textbf{Spain} and \textbf{Portugal} joined as new democracies---widened North--South gaps inside the Community and justified \textbf{Structural Funds} and later \textbf{Cohesion Fund} transfers (richer net contributors helping infrastructure and convergence in poorer regions).\\
\textbf{Single European Act (1986):} targeted a \textbf{``market without frontiers''} for goods/services/capital/people (four freedoms logic); extended \textbf{supranational} decision tools (more \textbf{QMV} in Council areas) and strengthened \textbf{European Parliament} participation relative to consultation-only infancy.\\
\textbf{Schengen Agreement:} signed \textbf{1985}, operational removal of internal border controls deepened across 1990s (\textbf{1995} milestone in many tellings); common external border and visa practice for participating states---see also \S 7 on membership breadth.\\
\textbf{UK counter-weight:} \textbf{Margaret Thatcher} embodied suspicion of ``\textbf{European superstate},'' prizing liberalized competition and national discretion over harmonized \textbf{social} standards---shadow of earlier Luxembourg-style vetoes in new guise.\\
\textbf{Treaty on European Union / \textbf{Maastricht} (signed \textbf{1992}, ratification drama after):} formally created the \textbf{European Union} name atop the Communities.\\
\textbf{Three pillars (classic Maastricht architecture):} (1) \textbf{European Communities} / economic integration including single market completion and \textbf{EMU} roadmap toward the \textbf{euro}; (2) \textbf{Common Foreign and Security Policy (CFSP)}; (3) \textbf{Justice and Home Affairs (JHA)} cooperation (later refined into today's AFSJ).\\
\textbf{EMU convergence:} \textbf{Maastricht criteria} on inflation, deficits/debt, interest rates, and exchange stability as euro-entry conditions (exam anchor).\\
\textbf{Opt-outs \& variable geometry:} \textbf{UK} and \textbf{Denmark} negotiated exemptions (e.g., \textbf{Social Charter} / Protocol opt-outs, single currency participation), institutionalizing a \textbf{multi-speed} Europe.\\
\textbf{Social Charter:} worker-rights principles contested by UK Conservatives during Maastricht---symbol of market vs social Europe divide.\\
\textbf{Cohesion strategy:} explicit \textbf{redistribution} rhetoric to keep enlargement politically sustainable when poorer members joined the common market.\\
\textbf{``Backdoor federalism'':} Delors-era strategy---deepening \textbf{economic} interdependence and shared rulebooks forces common legislation, approximating federal outcomes even when ``United States of Europe'' language is avoided.\\
\textbf{Exam phrasing:} ``1986--93 yokes market liberalization (SEA) to treaty constitutionalization (Maastricht): economics leads, while CFSP/JHA test how far political union follows.''
}

\cornell{14) Critical Perspectives: Vinen, Anderson, Bickerton, and the Democratic Deficit}
{How does Vinen reinterpret 1989 ``unity''?\\
What is Anderson's critique of elite-driven integration?\\
Why does the seminar spotlight the CJEU/ECJ?\\
What is Bickerton's ``memberization'' and EU scapegoating?\\
What socio-economic split does 1990s Czech data illustrate?\\
How does ``two faces'' explain Euroskepticism?}
{
\textbf{Counter-narrative goal:} move beyond celebratory ``official'' EU history toward structural, class, and legitimacy questions (readings such as \textbf{Richard Vinen}, \textbf{Perry Anderson}).\\
\textbf{Vinen on 1989:} fall of the \textbf{Berlin Wall} did not simply merge East and West; it generated \textbf{new cleavages}---\textbf{urban vs rural}, \textbf{wealthy vs marginalized}, and \textbf{generational} splits (youth adopting Western mobility/consumption vs older cohorts losing \textbf{socialist-era security} narratives).\\
\textbf{Illustrative data (1990s Czech Republic):} \textbf{Prague} unemployment reportedly ca.\ \textbf{0.32\%} amid global-city integration, versus national rate ca.\ \textbf{6.1\%}---spatial inequality inside one accession state.\\
\textbf{Anderson / democratic deficit:} integration advanced through \textbf{vanguard elites}, diplomats, and \textbf{Commission} technocrats plus \textbf{CJEU} jurisprudence more than mass democratic mandates; technical idioms (``economic cooperation'') allegedly \textbf{mask} deeper political choices voters might resist.\\
\textbf{CJEU power thesis:} often portrayed as uniquely consequential for enforcing \textbf{EU law primacy} over national statutes; critical scholars argue case law historically tilted toward \textbf{market-opening} (\textbf{neoliberal} integration) with uneven protection of \textbf{labor} or welfare (debate, not settled fact).\\
\textbf{Bickerton / ``memberization'':} national executives orient accountability toward \textbf{Brussels peer bargains} and package deals; domestic publics experience policy as external imposition; governments \textbf{scapegoat} ``Brussels made me do it'' for painful reforms (\textbf{institutional scapegoating}).\\
\textbf{Luxembourg shadow:} national \textbf{vital interest} vetoes (\S 10--11) coexist with judicial supranationalism---producing hybrid legitimacy tensions Anderson-style critics emphasize.\\
\textbf{Controversial genealogy tangent (as raised in seminar):} some discussions probe morally compromised mid-century technocratic networks among early architects---treat as historiographically disputed lecture provocation, not uncritically factual in exam answers.\\
\textbf{``Two faces'' conclusion:} for \textbf{urban Europeanized elites}, integration means mobility and opportunity; for \textbf{rural/marginalized} groups, removal of the Wall can coincide with \textbf{invisible borders} of class, age, and market access---fueling contemporary \textbf{Euroskepticism}.\\
\textbf{Exam phrasing:} ``Critical EU scholarship pairs legal integration (primacy, single market) with political sociology (who wins, who narrates, who pays), explaining gaps between institutional prose and democratic felt experience.''
}

\cornell{15) Gibraltar, Spain--UK Diplomacy, and the 1982 Border (Primary Sources)}
{What is Gibraltar geographically and legally?\\
Why did Franco close the border (1969)?\\
How did democratization change Spanish incentives?\\
Who reopened the border in 1982 and how fully?\\
What external crises framed 1982 diplomacy?\\
What are Lisbon (1980) and Brussels (1984) in the ``Brussels Process''?\\
How did student source groups split the evidence?\\
Why does the seminar treat integration as leverage?}
{
\textbf{Seminar frame (guest / jigsaw method):} archival seminar (e.g., guest professor Luz Martinez Kampo in transcript) where student clusters reconstruct \textbf{Spain--UK} bargaining over \textbf{Gibraltar} from \textbf{declassified} minutes, telegrams, and newspapers---complements \S 12 territorial notes.\\
\textbf{Geography:} British \textbf{Overseas Territory} (ca.\ 2.6 sq mi) on southern Iberian tip; linked to Spain by sandy \textbf{isthmus} hosting airport/border infrastructure.\\
\textbf{1969 closure:} dictator \textbf{Francisco Franco} sealed the frontier to pressure UK \textbf{sovereignty} concessions; humanitarian/economic shock---divided families, strangled Gibraltarian commerce.\\
\textbf{Post-1975 transition:} elected Spanish elites prioritized \textbf{European Community} accession and \textbf{NATO} alignment; London leveraged Madrid's need for Western/EC credibility to insist on progressive \textbf{border liberalization}.\\
\textbf{1982 reopening:} negotiated under Spanish socialist \textbf{Felipe González} and UK Premier \textbf{Margaret Thatcher}; initial step was \textbf{pedestrian-only} partial opening before fuller goods/vehicle regimes evolved.\\
\textbf{Falklands factor:} simultaneous \textbf{1982} war over South Atlantic islands raised temperature around colonial sovereignty disputes---parallel optics for UK resolve on overseas possessions.\\
\textbf{NATO synchronization:} \textbf{Spain joined NATO in 1982}, same calendar year as partial border opening---signals shared Western security basket even while sovereignty quarrels continued.\\
\textbf{Brussels Process anchors:} \textbf{Lisbon Agreement (1980)} and \textbf{Brussels Declaration (1984)} structured talks to air all differences (including ultimate sovereignty) while formally protecting \textbf{Gibraltarian} interests and \textbf{self-determination} language (UK emphasis).\\
\textbf{Source labs (student groupings):} \textit{Thatcher minutes}---UK insistence on resident rights and self-rule for \textbf{ca.\ 38{,}000} people; \textit{Spanish telegrams}---reconciling European vocation with historic \textbf{``Rock''} territorial claim; \textit{newspaper comparison}---Spanish press stress \textbf{national dignity}, British titles stress \textbf{rights/sovereignty} framing.\\
\textbf{Catalyst thesis:} \textbf{European integration} and reputational demands of ``good neighborhood'' within impending \textbf{single-market} family functioned as diplomatic \textbf{carrot/stick} to unwind a Cold War-era blockade incompatible with modernized Spanish identity.\\
\textbf{Exam phrasing:} ``Gibraltar shows imperial residue meeting Community ambition: sovereignty talk cannot be divorced from accession bargains, NATO politics, and local self-determination claims documented in archival evidence.''
}

\cornell{16) EU Historiography: Haas, Moravcsik, Gilbert, Anglosphere, and Brexit Ideology}
{What is the ``orthodox'' EU story and its critique?\\
How does Haas-style ``spillover'' explain integration?\\
How does Moravcsik-style structural/intergovernmental theory differ?\\
What is the Anglosphere / CANZUK imaginary post-Brexit?\\
How does Five Eyes relate?\\
What does Gilbert warn about teleological history?\\
What is the ``European train'' metaphor?\\
How did 2016--2020 US politics complicate UK strategy?}
{
\textbf{Historiography divide:} ``orthodox'' / Whig narrative presents integration as a progressive \textbf{march} toward peace and unity (heroes like \textbf{Monnet}/\textbf{Schuman}, villains/obstruction figures like \textbf{de Gaulle}/\textbf{Thatcher}); critical readers call this \textbf{teleological}---history written as if ``United States of Europe'' were predetermined.\\
\textbf{Haas / neofunctional spillover:} deepening in one sector (e.g., coal/steel) creates functional pressures and alliances that drag politics into adjacent sectors (customs union, single market logic); \textbf{critique:} over-deterministic, underplays national veto moments and electoral shocks.\\
\textbf{Moravcsik / materialist intergovernmental reading (termed ``structuralism'' in class):} treats grand bargains as outcomes of dominant domestic producer groups and executive \textbf{preference aggregation}---integration ``locks in'' market gains for exporters; missing cultural/identity variables.\\
\textbf{Gilbert caveat:} historiography should not assume a straight line; if integration ever unraveled (League/Soviet analogies in discussion), narratives might flip overnight to ``inevitable failure''---revealing presentism risk (\S 14 critical mood).\\
\textbf{``Stagnation'' decades:} 1970s--80s \textbf{Euro-pessimism} and patchy momentum complicate Whig triumphalism yet often vanish from textbook summaries.\\
\textbf{Anglosphere idea:} post-\textbf{Brexit} ideological package (\textbf{CANZUK}: Canada, Australia, New Zealand, UK) stressing English-language law, culture, intelligence ties as alternative to \textbf{supranational} EU obedience.\\
\textbf{Five Eyes (FVEY):} US--UK--Canada--Australia--NZ signals pact cited as template for deeper trade/security alignment fantasies beyond Brussels.\\
\textbf{Trump-era jolt (2016--2020):} ``America First'' volatility meant UK could not bank on stable Atlantic partnership while simultaneously exiting the EU \textbf{single market}---Anglosphere roadmap collided with unreliable Washington leadership in lecture framing.\\
\textbf{Train metaphor:} integration sold as a speeding \textbf{train} states must board or miss; critic asks who accepts perpetual acceleration toward an \textbf{unspecified terminus}---uncertainty feeding contemporary \textbf{Euroskepticism}.\\
\textbf{Sovereignty punchline:} Brexit discourse elevated reclaimed \textbf{sovereignty} versus ``forced compliance'' with EU law---mirrors Anglosphere pitch.\\
\textbf{Exam phrasing:} ``Competing IR theories and historiographies interpret the same treaties differently: spillover vs intergovernmental bargaining vs post-imperial English-speaking alliances each carry distinct policy prescriptions.''
}

\cornell{17) Modern EU (1993--2026): Enlargement East, Crises, Euro, and Turkey}
{What enlargements followed Maastricht?\\
Which institutional shocks punctuated the 1990s--2000s?\\
How did the euro and ECB reshape crisis politics?\\
What are the main arguments for and against Turkish accession?\\
Why is Cyprus a structural veto?\\
How does ``Copenhagen'' frame accession today?\\
What is the ``geopolitical bouncer'' idea?}
{
\textbf{Post-Maastricht widening:} \textbf{1995} accession of \textbf{Austria, Finland, Sweden} (former neutrals pivoting into single-market rules); \textbf{2004 ``Big Bang''}---ten states, largely ex-communist Central/Eastern Europe and Mediterranean entrants---shifted geopolitical center of gravity \textbf{eastward}.\\
\textbf{Commission crisis (1999):} entire \textbf{Santer Commission} resigned amid fraud/mismanagement allegations---early transparency and oversight stress test.\\
\textbf{Constitutional moment failure:} \textbf{2005} French and Dutch referendums defeated the draft \textbf{European Constitution}; negotiators repackaged core reforms as the \textbf{Lisbon Treaty (2007)} ``Plan B'' (\S 13).\\
\textbf{Euro launch:} banknotes/coins in \textbf{2002}; \textbf{Eurozone} counts \textbf{20} member states using the single currency in the 2026 framing cited in class (number evolves).\\
\textbf{ECB / Frankfurt technocracy:} \textbf{European Central Bank} exercises powerful monetary instruments; \textbf{Greek debt crisis} era showcased conditionality and tension between democratic budgeting and ECB--Troika oversight (\textbf{technocracy} critique---detail from course readings).\\
\textbf{Copenhagen criteria (contemporary accession bar):} proves a ``fully functioning democracy''---(1) stable institutions and rule of law; (2) human rights including minorities and speech; (3) market economy able to withstand competitive pressure; full adoption of \textbf{acquis communautaire} remains the legal hurdle (\S 7--8 echo).\\
\textbf{Turkey---strategic ``pros'':} bridge state toward Middle East/Central Asia, energy corridor leverage; large \textbf{NATO} army; youthful demographics potentially offsetting EU aging.\\
\textbf{Turkey---``cons'':} democratic backsliding under \textbf{Erdoğan}, political prisoners, judicial independence doubts; geography mostly \textbf{Asian} with unstable southern borders (Syria/Iraq); identity debates contrasting Islamic heritage with EU ``Christian/Enlightenment'' self-image in polemical discourse.\\
\textbf{Kemalist backdrop:} \textbf{secularism} (Kemalism) as founding republican principle now contested by religious-conservative politics---affects rule-of-law scoring in accession talk.\\
\textbf{Cyprus deadlock:} internationally recognized \textbf{Republic of Cyprus} (Greek-majority south) is EU member while north \textbf{TRNC} (Turkish-backed) lacks recognition; stalemate yields Cypriot/Greek \textbf{veto} leverage against chapters touching Turkey.\\
\textbf{``Geopolitical bouncer'':} Turkey (lecture also mentions Morocco/Libya parallels) manages migration and security \textbf{externalization} for the EU---like Mexico--US analogies---exchanging stabilization work for funds/concessions despite stalled membership.\\
\textbf{QMV reminder:} Lisbon-era Council rules widen majority voting domains, yet politically sensitive dossiers (enlargement, foreign policy) still flirt with unanimity traps.\\
\textbf{Exam phrasing:} ``Post-2004 EU juggles normative enlargement promises with hard-border politics: Turkey lays bare the clash between Copenhagen ideals, identity politics, and neighborhood stabilization deals.''
}

\cornell{18) Early 2026 Lens: Multilateral Stress, Energy War, Kelsen vs Schmitt}
{Which milestones anchor the ``modern EU'' story in seminar recap?\\
What is fraying about the post-1945 rules-based order?\\
How do von der Leyen and Council leaders differ in diagnosis?\\
What did COVID reveal about EU power and vulnerability?\\
How did Ukraine 2022 reshape energy geopolitics?\\
Why invoke Kelsen and Schmitt today?\\
What are the emerging spheres-of-influence blocs?\\
What is the ``Gilded Age / Progressive Era'' analogy?\\
What does ``end of the tree'' claim about war and law?}
{
\textbf{Milestone refresh:} \textbf{1993} \textbf{EU} treaty formalization (\textbf{Maastricht}); \textbf{1999} \textbf{EMU} Stage Three and \textbf{ECB} operations (euro banknotes in \S 17 usually dated \textbf{2002}); \textbf{2004} eastern \textbf{Big Bang}; \textbf{2007} \textbf{Lisbon} ``Plan B'' after \textbf{2005} constitutional referenda.\\
\textbf{Rules-based order under strain:} seminar asks whether \textbf{UN}/\textbf{WTO}/\textbf{EU} legalism survives when great powers treat law as optional---\textbf{multilateralism} versus \textbf{unilateralism} and competing \textbf{sovereignty} claims (early \textbf{2026} discussions in transcript).\\
\textbf{Von der Leyen inflection:} Commission discourse acknowledges partners (notably US) drifting from treaty fidelity; risk that Europe remains uniquely \textbf{rule-bound} while rivals exercise raw power.\\
\textbf{Costa / member-government counter-voice:} premiers (e.g., Portuguese leadership cited) stress budgetary solidarity and political feasibility---\textbf{Council} intergovernmentalism tempers federal appetite when crises hit national welfare systems.\\
\textbf{Realist overlay:} third strand treats integration rhetoric as luxury when \textbf{hard-power blocs} redraw maps regardless of legal order ideals (\S 16 historiography complements).\\
\textbf{COVID moment:} joint vaccine procurement and \textbf{recovery funds} showed unprecedented fiscal centralization yet exposed \textbf{tourism} and \textbf{supply-chain} global dependency.\\
\textbf{Energy after 2022:} full-scale invasion of \textbf{Ukraine} forced rapid \textbf{decoupling} from \textbf{Russian} pipeline gas; \textbf{Spain} spotlighted for \textbf{renewables} capacity while entire Union remains price-taker on \textbf{liquefied gas} and \textbf{global} markets---\textbf{strait/canal choke points} (Suez, Gibraltar, Panama) remembered as trade risk nodes.\\
\textbf{Kelsen vs Schmitt:} \textbf{Kelsen}---normative hierarchy and universal rules; \textbf{Schmitt}---\textbf{sovereign decides the exception}; classroom sense that crises (pandemic, war, migration surges) tilt politics toward \textbf{Schmittian} executive discretion ``over'' law.\\
\textbf{Neo-spheres sketch:} US \textbf{America First} trade weaponization; \textbf{China} Indo-Pacific/Africa leverage; \textbf{Russia} kinetic pressure on EU's eastern rim---19th-century-style \textbf{spheres of influence} vocabulary returns.\\
\textbf{US history analogy:} \textbf{Gilded Age} (rapid growth, weak center, corruption optics) vs \textbf{Progressive Era} (activist federal repair)---students ask whether EU is stuck in bureaucratic \textbf{Gilded} softness or edging toward centralized \textbf{Progressive} emergency politics.\\
\textbf{``End of the tree'' motif:} reformist student metaphor---official teleology of ever greener European progress ``hits the ground'' when simultaneous wars (Ukraine, Middle East) show great powers unconstrained by universal law; choice framed starkly: EU must become a \textbf{geopolitical power} or risk marginality as mere \textbf{legal entity}.\\
\textbf{Macro echo:} \textbf{stagflation} fears in \textbf{2020s} linked rhetorically to \textbf{1973} supply shocks (\S 11--12).\\
\textbf{Exam phrasing:} ``Contemporary EU theory must pair lawyers' integration scripts with energy and security shocks: the question is whether law can still constitute power or only chase it.''
}

\cornell{19) European Parliament, Co-Decision, and Democratic Deficit}
{What is the democratic deficit in EU studies?\\
How did the EP move from ``talking shop'' to co-legislator?\\
Why are political groups ideological, not national?\\
How does the institutional triangle produce EU law?\\
What was the housing-crisis group exercise?\\
Which EP ``families'' anchor majorities?\\
How do subsidiarity and acquis relate?\\
Why does an ``education gap'' matter for legitimacy?}
{
\textbf{Democratic deficit:} perceived distance between opaque Brussels procedures and citizens' sense of control; historically tied to dominance of appointed \textbf{Commission} executives and \textbf{Council} diplomats versus weak \textbf{parliamentary} scrutiny (nuanced after Lisbon---jury still debated; links to \S 14 Anderson-style critique).\\
\textbf{EP timeline:} \textbf{1952}---Assembly largely appointed consultative chamber without blocking law; \textbf{1979}---first \textbf{direct elections} every five years create EU-wide mandate claims; \textbf{Maastricht} onward and especially \textbf{Lisbon} consolidate \textbf{Ordinary Legislative Procedure (OLP)} / \textbf{co-decision}: Commission \textbf{proposes}, Parliament and Council must \textbf{both} adopt text for most single-market and social dossiers.\\
\textbf{Political groups:} MEPs sit in transnational \textbf{ideological} clubs (e.g., Spanish and German socialists in \textbf{S\&D}) rather than national delegations---negotiate compromises and committee leadership inside that logic.\\
\textbf{Institutional triangle (lawmaking cartoon):} (1) \textbf{Commission}---exclusive right of initiative / quasi-executive; (2) \textbf{Council}---member-state governments defend national bargains; (3) \textbf{Parliament}---represents \textbf{citizens} as directly elected check; balances mirror federalism debates without a single EU president matching US tripartite myth.\\
\textbf{Major group map (lecture snapshot):} \textbf{EPP}---center-right Christian democrats/conservatives; \textbf{S\&D}---center-left social democrats; \textbf{Renew Europe}---liberal centrists; \textbf{Greens/EFA}---environmentalists and regionalists; \textbf{ECR}---right-wing sovereigntist reformers; \textbf{Patriots for Europe}---far-right nationalist family (new constellation in early-2026 classroom talk); \textbf{The Left}---radical left including democratic socialists/communists.\\
\textbf{Housing crisis exercise:} student teams compared party-group platforms on affordability---\textbf{Greens}, \textbf{EPP}, \textbf{S\&D}, etc.---with \textbf{Vienna}-style high-share \textbf{public/social housing} (ca.\ \textbf{50\%} share in transcript's benchmark) as reference for ``\textbf{Social Europe}'' ambitions vs market-led supply solutions.\\
\textbf{Subsidiarity:} EU acts only when objectives cannot be sufficiently achieved at lower tiers---key rhetorical limit on federal overreach in housing and welfare debates.\\
\textbf{Acquis hook:} once legislation passes OLP it joins enforceable \textbf{acquis} monitored by Commission and \textbf{CJEU}.\\
\textbf{Strategic insight---``education gap'':} even after empowerment, citizens often cannot trace \textbf{OLP} steps; low media literacy on ``who decides'' amplifies \textbf{deficit feelings} and populist blame games despite real EP powers.\\
\textbf{Exam phrasing:} ``Parliament's rise proves formal democracy can deepen, yet legitimacy still depends on public knowledge of co-decision mechanics---institutional reform without civic pedagogy leaves the deficit narrative alive.''
}

\cornell{20) CJEU, Primacy, Backsliding (Hungary/Poland), and Rule-of-Law Tools}
{What is the CJEU's core mandate vis-\`a-vis 27 legal orders?\\
Which founding doctrines anchor EU legal autonomy?\\
How do preliminary references knit national and EU courts?\\
What is democratic backsliding in the EU context?\\
What is Article 7 and why is it ``nuclear''?\\
How does rule-of-law conditionality differ?\\
How does Gros's ``double standard'' critique work?\\
What are EU--US high-court parallels?}
{
\textbf{CJEU function:} \textbf{Court of Justice of the European Union} secures uniform interpretation/application of \textbf{EU law} across \textbf{27} members---judicial \textbf{gap-filling} when treaties/regulations leave ambiguities (\S 2 baseline).\\
\textbf{Primacy (\textit{Costa v.\ ENEL}, 1964):} conflicts between valid EU norms and later national statutes resolve in favor of \textbf{EU law}---constitutional tension with some national courts remains a live issue.\\
\textbf{Direct effect (\textit{Van Gend en Loos}, 1963):} certain treaty provisions create \textbf{subjective rights} plaintiffs can invoke before domestic judges---underpins decentralized enforcement.\\
\textbf{Preliminary ruling procedure (Art.\ 267 TFEU):} national courts \textbf{refer} questions of interpretation/validity to Luxembourg---creates cooperative ``bridge'' likened to ongoing \textbf{juridification} of politics.\\
\textbf{Democratic backsliding:} elected executives erode checks---capturing courts, media, civil society---exemplified in seminar by \textbf{Fidesz}-led \textbf{Hungary} and (in the lecture timeline) \textbf{PiS}-era policies in \textbf{Poland} even as Warsaw governments rotate.\\
\textbf{``Illiberal democracy'' tag:} retains electoral shells while constraining pluralism and rights---feeds friction with Article 2 values.\\
\textbf{Article 7 TEU (``nuclear option''):} procedure to determine ``serious and persistent breach'' of values and ultimately \textbf{suspend} Council \textbf{voting rights}---rarely used because it demands extraordinary political consensus among peers (\textbf{unanimity} for key steps in standard teaching).\\
\textbf{Rule-of-law conditionality regulation:} \textbf{budget-protection} tool allowing Commission/Council to \textbf{freeze cohesion/recovery funds} when breaches threaten EU financial management---material leverage distinct from Article 7 symbolism.\\
\textbf{Gros / inconsistency critique:} EU appears swift/harsh on \textbf{economic governance} (e.g., Greek-style crisis discipline) yet hesitant on \textbf{values enforcement}---\textbf{double standard} allegation feeding populist \textbf{sovereignty} narratives in Budapest/Warsaw.\\
\textbf{QMV reminder (Council):} many enforcement dossiers now use weighted majorities (lecture shorthand: \textbf{55\%} of states representing \textbf{65\%} of EU population) rather than unanimity---contrast Article 7 politics.\\
\textbf{CJEU vs US Supreme Court analogy:} both apex courts legitimate integration/federal bargains; differences include CJEU's comparative \textbf{opacity} (no published dissenting opinions in classic mode) and obligation to harmonize \textbf{27} distinct constitutional orders vs a single federation.\\
\textbf{Legal certainty \& sovereignty:} capitals invoke domestic \textbf{sovereignty} against Brussels while investors demand predictable EU \textbf{legal certainty}---backsliding threatens both market confidence and pluralist democracy.\\
\textbf{Exam phrasing:} ``Integration through law now means policing judiciaries: the CJEU defends the acquis, yet Article 7 stalemates force the Union toward budgetary coercion when values and money intersect.''
}

\section*{Summary Block (Cornell Bottom)}
\begin{itemize}
  \item High-scoring answers integrate institutions, law, and political economy in one argument line.
  \item Use formula: concept definition -> mechanism -> example -> implication.
  \item Allocate extra revision time here (focus-most course group).
  \item Finals tile status: no course-specific image tiles detected in current upload batch; keep this file aligned with class notes/readings.
\end{itemize}

\end{document}
